
\documentclass[11pt]{article}
\usepackage[margin=1in]{geometry}
\usepackage{amsmath,amssymb,amsthm,mathtools,booktabs,array,enumitem,hyperref,microtype}
\usepackage[T1]{fontenc}
\usepackage{lmodern}

\newtheorem{theorem}{Theorem}[section]
\newtheorem{proposition}[theorem]{Proposition}
\newtheorem{lemma}[theorem]{Lemma}
\newtheorem{corollary}[theorem]{Corollary}
\newtheorem{definition}[theorem]{Definition}
\newtheorem{remark}[theorem]{Remark}

\newcommand{\Z}{\mathbb Z}
\newcommand{\Q}{\mathbb Q}
\newcommand{\C}{\mathbb C}
\newcommand{\N}{\mathbb N}
\newcommand{\Zi}{\Z[i]}
\newcommand{\Qi}{\Q(i)}
\newcommand{\HA}{\mathcal H_A}
\newcommand{\HR}{\mathcal H}
\newcommand{\CA}{\mathcal C_A}
\newcommand{\Val}{\operatorname{Val}_A}
\newcommand{\Can}{\operatorname{Can}_A}
\newcommand{\Real}{\operatorname{Real}}
\newcommand{\coeff}{\operatorname{coeff}}

\title{\textbf{A Hybrid CNRS-A/CNRS-H Representation Theorem}\\
\large Canonical Complex Coefficients and the Hurwitz-Series Analytic Carrier}
\author{Donald G. Palmer}
\date{July 2026}

\begin{document}
\maketitle

\begin{abstract}
CNRS-A and CNRS-H perform different mathematical roles.  CNRS-A supplies canonical positional representations of complex coefficients, while CNRS-H supplies the exponential-generating-function index structure in which differentiation is a coefficient shift.  We formalize this as a two-carrier architecture.  Given a coefficient ring $R$ and a canonical CNRS-A representation set $\CA$ equipped with a bijective value map $\Val:\CA\to R$, we define the hybrid space $\HA(\CA)=\CA^{\N}$.  Transporting the Hurwitz product and differential operators through $\Val$ gives $\HA(\CA)$ a unique differential-algebra structure for which the coefficientwise value map is an isomorphism
\[
\HA(\CA)\cong\HR(R).
\]
Every formal CNRS-H series over $R$ therefore has a unique representation as a sequence of canonical CNRS-A coefficients.  The theorem separates arithmetic representation from analytic order, proves compatibility with addition, Hurwitz multiplication, differentiation, integration, units, composition, and formal inversion, and identifies the conditions under which coefficientwise completeness transfers to the hybrid system.  An optional analytic realization into complex functions is treated separately and requires a chosen embedding and growth conditions.  The result establishes the mathematically natural hybrid architecture without requiring a literal base-$e$ positional numeral system.
\end{abstract}

\section{Motivation}

The CNRS programme contains two structures that should not be conflated.

\begin{enumerate}[label=(\roman*)]
\item CNRS-A is a positional representation layer for complex arithmetic, based in the present implementation on the Gaussian base
\[
\beta=-2+i.
\]
It supplies canonical finite, Laurent, or canonical eventually periodic representations for the coefficient domain under consideration.

\item CNRS-H is a formal exponential-generating-function layer.  Its elements are coefficient sequences
\[
(a_0,a_1,a_2,\ldots),
\]
interpreted formally as
\[
\sum_{n\ge0}a_n\frac{\rho^n}{n!}.
\]
Its derivative is the left shift.
\end{enumerate}

The correct unification is therefore not necessarily a single positional base doing both jobs.  It is a two-carrier representation in which CNRS-A carries coefficient values and CNRS-H carries analytic order.

\section{Coefficient Representation Data}

Let $R$ be a commutative ring with identity.

\begin{definition}[Canonical CNRS-A coefficient system]
A canonical CNRS-A coefficient system over $R$ is a triple
\[
(\CA,\Val,\Can)
\]
such that:
\begin{enumerate}[label=(\alph*)]
\item $\CA$ is a set of canonical coefficient representations;
\item $\Val:\CA\to R$ is a bijection;
\item $\Can:R\to\CA$ is its inverse.
\end{enumerate}
\end{definition}

Thus
\[
\Val(\Can(r))=r,\qquad
\Can(\Val(c))=c.
\]

Examples of possible coefficient rings include:
\[
R=\Zi,\qquad R=\Qi,\qquad R=\mathcal O_\beta,\qquad R=K_\beta,
\]
provided the chosen CNRS-A representation layer is canonical and bijective on that domain.

\begin{remark}
The theorem below is representation-theoretic.  It does not assume that every complex number in $\C$ already has a canonical CNRS-A representation.  It applies to any coefficient ring for which such a canonical representation map has been established.
\end{remark}

Transport the ring operations to $\CA$ by
\[
c\oplus_A d
=
\Can(\Val(c)+\Val(d)),
\]
\[
c\otimes_A d
=
\Can(\Val(c)\Val(d)),
\]
and
\[
\ominus_A c
=
\Can(-\Val(c)).
\]

\begin{proposition}
With these transported operations, $\CA$ is a commutative ring and $\Val:\CA\to R$ is a ring isomorphism.
\end{proposition}

\begin{proof}
All ring axioms follow by transport across the bijection $\Val$.  For example,
\[
\Val((c\oplus_A d)\oplus_A e)
=
(\Val(c)+\Val(d))+\Val(e),
\]
which equals
\[
\Val(c)+(\Val(d)+\Val(e))
=
\Val(c\oplus_A(d\oplus_A e)).
\]
Injectivity of $\Val$ yields associativity.  The remaining axioms are identical.
\end{proof}

\section{The Hurwitz-Series Carrier}

Define
\[
\HR(R)=R^\N.
\]

Addition is coefficientwise:
\[
(a+b)_n=a_n+b_n.
\]

The Hurwitz product is
\[
(a*b)_n
=
\sum_{k=0}^n
\binom{n}{k}a_kb_{n-k}.
\]

The derivative and integral with chosen constant $c\in R$ are
\[
D(a_0,a_1,a_2,\ldots)
=
(a_1,a_2,a_3,\ldots),
\]
\[
I_c(a_0,a_1,a_2,\ldots)
=
(c,a_0,a_1,\ldots).
\]

When $R$ is a $\Q$-algebra, the formal exponential-generating-function realization is
\[
\mathcal E_R(a)
=
\sum_{n\ge0}a_n\frac{\rho^n}{n!}.
\]

The prior CNRS-H algebra theorem establishes that $\HR(R)$ is a commutative differential algebra under these operations.

\section{The Hybrid Representation Space}

\begin{definition}
The hybrid CNRS-A/CNRS-H representation space is
\[
\HA(\CA)=\CA^\N.
\]
An element
\[
\mathbf c=(c_0,c_1,c_2,\ldots)
\]
is called a hybrid coefficient string.
\end{definition}

Define the coefficientwise value map
\[
\Phi:\HA(\CA)\to\HR(R)
\]
by
\[
\Phi(\mathbf c)
=
(\Val(c_0),\Val(c_1),\Val(c_2),\ldots).
\]

Its inverse is
\[
\Psi:\HR(R)\to\HA(\CA),
\qquad
\Psi((a_n)_{n\ge0})
=
(\Can(a_n))_{n\ge0}.
\]

\begin{theorem}[Hybrid representation theorem]
The maps $\Phi$ and $\Psi$ are mutually inverse bijections.  Consequently every Hurwitz series over $R$ has a unique hybrid representation as a sequence of canonical CNRS-A coefficients:
\[
(a_n)_{n\ge0}
\longleftrightarrow
(\Can(a_n))_{n\ge0}.
\]
\end{theorem}

\begin{proof}
For $\mathbf c=(c_n)$,
\[
\Psi(\Phi(\mathbf c))
=
(\Can(\Val(c_n)))_{n\ge0}
=
(c_n)_{n\ge0}.
\]
For $a=(a_n)$,
\[
\Phi(\Psi(a))
=
(\Val(\Can(a_n)))_{n\ge0}
=
(a_n)_{n\ge0}.
\]
Thus the maps are inverse bijections.
\end{proof}

\section{Transported Differential-Algebra Structure}

Define hybrid addition by
\[
(\mathbf c\oplus_H\mathbf d)_n
=
c_n\oplus_A d_n.
\]

Define the hybrid Hurwitz product by
\[
(\mathbf c\star_H\mathbf d)_n
=
\bigoplus_{k=0}^{n}
\Can\!\left(
\binom{n}{k}
\right)
\otimes_A c_k\otimes_A d_{n-k},
\]
where the finite sum is evaluated in the transported coefficient ring.

Define
\[
D_H(c_0,c_1,c_2,\ldots)
=
(c_1,c_2,c_3,\ldots),
\]
and
\[
I_{H,\gamma}(c_0,c_1,c_2,\ldots)
=
(\gamma,c_0,c_1,c_2,\ldots)
\]
for $\gamma\in\CA$.

\begin{theorem}[Differential-algebra isomorphism]
There is a unique differential-algebra structure on $\HA(\CA)$ for which
\[
\Phi:\HA(\CA)\to\HR(R)
\]
is an isomorphism.  It is precisely the structure defined above.  In particular,
\[
\Phi(\mathbf c\oplus_H\mathbf d)
=
\Phi(\mathbf c)+\Phi(\mathbf d),
\]
\[
\Phi(\mathbf c\star_H\mathbf d)
=
\Phi(\mathbf c)*\Phi(\mathbf d),
\]
and
\[
\Phi(D_H\mathbf c)=D\Phi(\mathbf c).
\]
\end{theorem}

\begin{proof}
For addition the claim is coefficientwise.  For multiplication,
\[
\Val((\mathbf c\star_H\mathbf d)_n)
=
\sum_{k=0}^{n}
\binom{n}{k}
\Val(c_k)\Val(d_{n-k}),
\]
which is the $n$th Hurwitz-product coefficient.  Differentiation is immediate from the shift.  Uniqueness follows because every operation making $\Phi$ an isomorphism must be the pullback of the corresponding operation on $\HR(R)$.
\end{proof}

\begin{corollary}[Leibniz rule]
\[
D_H(\mathbf c\star_H\mathbf d)
=
(D_H\mathbf c)\star_H\mathbf d
+
\mathbf c\star_H(D_H\mathbf d).
\]
\end{corollary}

\begin{corollary}[Integration]
\[
D_H I_{H,\gamma}(\mathbf c)=\mathbf c.
\]
If $\gamma$ is the canonical zero coefficient, then
\[
I_{H,0}D_H(\mathbf c)
=
\mathbf c-(c_0,0,0,\ldots).
\]
\end{corollary}

\section{Units, Composition, and Inversion}

\begin{proposition}[Multiplicative units]
A hybrid series $\mathbf c=(c_n)$ is invertible under $\star_H$ if and only if $\Val(c_0)$ is a unit in $R$.  Its inverse is unique and is obtained by transporting the Hurwitz-series inverse recurrence.
\end{proposition}

\begin{proof}
Apply the corresponding unit criterion in $\HR(R)$ through the isomorphism $\Phi$.
\end{proof}

Suppose now that $R$ is a $\Q$-algebra, or more generally that the integer and Bell-polynomial coefficients required below act meaningfully in $R$.

For $a,b\in\HR(R)$ with $b_0=0$, formal composition is defined by
\[
\mathcal E_R(a\circ b)
=
\mathcal E_R(a)\bigl(\mathcal E_R(b)\bigr)
\]
formally, equivalently through exponential Bell polynomials.

Define hybrid composition by transport:
\[
\mathbf c\circ_H\mathbf d
=
\Psi\bigl(\Phi(\mathbf c)\circ\Phi(\mathbf d)\bigr).
\]

\begin{theorem}[Composition compatibility]
Whenever formal composition is defined,
\[
\Phi(\mathbf c\circ_H\mathbf d)
=
\Phi(\mathbf c)\circ\Phi(\mathbf d).
\]
Composition is associative on its natural domain, and the chain rule holds:
\[
D_H(\mathbf c\circ_H\mathbf d)
=
(D_H\mathbf c\circ_H\mathbf d)\star_H D_H\mathbf d.
\]
\end{theorem}

\begin{proof}
The first statement is the definition.  Associativity and the chain rule follow from the corresponding results in $\HR(R)$ and injectivity of $\Phi$.
\end{proof}

\begin{theorem}[Compositional inverse]
Let $\mathbf g=(g_n)$ satisfy
\[
\Val(g_0)=0,
\qquad
\Val(g_1)\in R^\times.
\]
Then there exists a unique hybrid series $\mathbf h$ such that
\[
\mathbf g\circ_H\mathbf h
=
\mathbf h\circ_H\mathbf g
=
\boldsymbol\rho,
\]
where
\[
\boldsymbol\rho=(\Can(0),\Can(1),\Can(0),\ldots).
\]
\end{theorem}

\begin{proof}
Transport the formal compositional inverse theorem from $\HR(R)$.
\end{proof}

\section{Exponential Eigenfunctions}

For $\alpha\in R$, define
\[
E_\alpha=(1,\alpha,\alpha^2,\ldots)\in\HR(R),
\]
and its hybrid representation
\[
\mathbf E_\alpha^A
=
(\Can(1),\Can(\alpha),\Can(\alpha^2),\ldots).
\]

\begin{theorem}
The hybrid exponential eigenfunction satisfies
\[
D_H\mathbf E_\alpha^A
=
\Can(\alpha)\star_H\mathbf E_\alpha^A,
\]
where $\Can(\alpha)$ on the right is identified with the constant hybrid series
\[
(\Can(\alpha),\Can(0),\Can(0),\ldots).
\]
Moreover,
\[
\mathbf E_\alpha^A\star_H\mathbf E_\beta^A
=
\mathbf E_{\alpha+\beta}^A.
\]
\end{theorem}

\begin{proof}
Apply $\Phi$ and use the identities
\[
D E_\alpha=\alpha E_\alpha,
\qquad
E_\alpha*E_\beta=E_{\alpha+\beta}
\]
in the Hurwitz algebra.
\end{proof}

This theorem gives the exact algebraic bridge between CNRS-A coefficient arithmetic and the exponential structure central to CNRS-H.

\section{Topological Transfer}

Let $R$ be a complete metrizable topological ring, and suppose $\CA$ is equipped with the transported metric
\[
d_A(c,d)=d_R(\Val(c),\Val(d)).
\]

Define on $\HA(\CA)$
\[
d_{\HA}(\mathbf c,\mathbf d)
=
\sum_{n=0}^{\infty}
2^{-n-1}
\min\{1,d_A(c_n,d_n)\}.
\]

\begin{theorem}[Hybrid completeness]
If $R$ is complete, then $\HA(\CA)$ is complete.  The map
\[
\Phi:\HA(\CA)\to\HR(R)
\]
is an isometric homeomorphism for the corresponding product metrics.
\end{theorem}

\begin{proof}
The equality
\[
d_{\HA}(\mathbf c,\mathbf d)
=
d_{\HR}(\Phi(\mathbf c),\Phi(\mathbf d))
\]
is immediate from the definition of $d_A$.  Completeness follows from the coefficientwise completeness theorem for $\HR(R)$.
\end{proof}

For example, if
\[
R=\mathcal O_\beta\cong\Z_5
\]
or
\[
R=K_\beta\cong\Q_5,
\]
then the corresponding hybrid CNRS-A/CNRS-H formal-series space is complete.

\section{Analytic Realization}

The hybrid theorem is formal and representation-theoretic.  To obtain an ordinary complex function, additional data are required.

Let
\[
\iota:R\to\C
\]
be a chosen coefficient embedding or evaluation map.  Define, where convergent,
\[
\Real_\iota(\mathbf c)(\rho)
=
\sum_{n=0}^{\infty}
\iota(\Val(c_n))
\frac{\rho^n}{n!}.
\]

\begin{proposition}[Uniqueness of analytic realization]
On any domain where the series converges, the analytic function is uniquely determined by the hybrid representation.
\end{proposition}

\begin{proof}
The canonical coefficient sequence is unique, and an exponential power series is determined by its coefficients.
\end{proof}

However, existence of a nonzero radius of convergence requires growth control, for example
\[
\limsup_{n\to\infty}
\left(
\frac{|\iota(\Val(c_n))|}{n!}
\right)^{1/n}
<\infty.
\]

\begin{remark}
For $R=\Q(i)$, the usual embedding into $\C$ is available.  For a $\beta$-adic completion such as $K_\beta\cong\Q_5$, there is no continuous field embedding into $\C$ compatible with the respective topologies.  A hybrid object over $K_\beta$ is therefore naturally a non-Archimedean formal or analytic object unless additional, noncanonical coefficient interpretation is supplied.
\end{remark}

\section{Two-Carrier Uniqueness}

\begin{theorem}[Two-carrier uniqueness]
Fix:
\begin{enumerate}[label=(\roman*)]
\item a coefficient ring $R$;
\item a canonical CNRS-A coefficient system $(\CA,\Val,\Can)$;
\item the EGF basis $\{\rho^n/n!\}_{n\ge0}$.
\end{enumerate}
Then every element of $\HR(R)$ has exactly one hybrid representation
\[
\sum_{n\ge0}
\Val(c_n)\frac{\rho^n}{n!},
\qquad c_n\in\CA.
\]
If two hybrid strings determine the same formal series, then they agree coefficientwise.
\end{theorem}

\begin{proof}
Equality of formal exponential series implies equality of all coefficients.  Bijectivity of $\Val$ then implies equality of the canonical coefficient representations.
\end{proof}

This is the central representation theorem: uniqueness comes from canonical coefficient normalization together with a fixed analytic basis.

\section{What the Theorem Does and Does Not Establish}

The theorem establishes:
\begin{enumerate}
\item a unique CNRS-A representation for every coefficient in the selected ring;
\item a unique hybrid representation for every Hurwitz series over that ring;
\item exact transport of arithmetic and differential operations;
\item coefficientwise completeness whenever the coefficient ring is complete;
\item a precise route to analytic realization when an embedding and growth conditions are supplied.
\end{enumerate}

It does not establish:
\begin{enumerate}
\item that all of $\C$ already has a single chosen canonical CNRS-A coefficient representation;
\item that a literal base-$e$ positional system is required;
\item that every formal CNRS-H object converges as an ordinary complex function;
\item that the $\beta$-adic and ordinary complex topologies coincide;
\item a physical interpretation for the hybrid representation.
\end{enumerate}

\section{Implications for CNRS}

The result supports the architectural statement
\[
\boxed{
\text{CNRS}
=
\text{CNRS-A coefficient carrier}
+
\text{CNRS-H analytic carrier}.
}
\]

The plus sign denotes structured composition, not identification.  CNRS-A solves the coefficient-representation problem.  CNRS-H solves the analytic-order and differential-shift problem.  Their combination is canonical once the coefficient domain and basis convention are fixed.

For implementation, the theorem recommends:
\begin{enumerate}
\item store each CNRS-H coefficient as a canonical CNRS-A object;
\item canonicalize after each coefficient arithmetic operation;
\item implement Hurwitz multiplication coefficientwise through CNRS-A arithmetic;
\item retain a strict distinction between formal series and analytic evaluation;
\item attach topology metadata to the coefficient domain.
\end{enumerate}

\section{Conclusion}

CNRS-A and CNRS-H need not be collapsed into one positional base.  The mathematically natural unification is a two-carrier representation theorem.  Canonical CNRS-A strings represent coefficient values; the Hurwitz-series basis represents analytic order.  The coefficientwise value map is a differential-algebra isomorphism, and completeness transfers from the coefficient ring to the hybrid formal-series space.  Ordinary complex-function realization is an additional layer governed by embedding and growth conditions.  This closes the formal architecture of the arithmetic/analytic bridge and provides a precise foundation for the corresponding Toolkit layer.

\section*{AI-Assistance Declaration}
AI tools were used to assist with theorem organization, derivation checking, drafting, and computational verification.  The author remains responsible for the mathematical claims, interpretation, and publication decisions.

\begin{thebibliography}{9}
\bibitem{Hurwitz}
A. Hurwitz,
standard foundational work on formal series and coefficient calculus.

\bibitem{Roman}
S. Roman,
\emph{The Umbral Calculus},
Academic Press.

\bibitem{KataiSzabo}
I. K\'atai and J. Szab\'o,
Canonical number systems for complex integers,
\emph{Acta Scientiarum Mathematicarum} 37 (1975), 255--260.

\bibitem{PalmerH}
D. G. Palmer,
Formal CNRS-H algebra,
CNRS working theorem paper, 2026.

\bibitem{PalmerCompleteness}
D. G. Palmer,
Metric and topological completeness in the CNRS architecture,
CNRS working theorem paper, 2026.

\bibitem{PalmerNormalization}
D. G. Palmer,
Canonical periodic normalization for CNRS,
CNRS working theorem paper, 2026.
\end{thebibliography}

\end{document}
